\paragraph{The four exploration invariants.}\label{ta:search-invariants}
For the partial table $\mathcal D$, queried-empty buckets are recorded by candidate consumption even though they are absent from $\operatorname{dom}\mathcal D$. Each stored nonempty bucket is exact and contains all modeled outcomes. We state invariants at loop boundaries, after discovery queues and the finite reverse-edge worklist have been updated.
\begin{enumerate}
\item \textbf{Uncontrollable completeness.} On first expansion of $q$, all uncontrollable candidates are consumed before promotion can use the state. No candidate is added later: the finite event set is determined by $q$. Therefore every $q\in X$ has its complete uncontrollable successor set known. The transient insertion into $X$ inside the expansion block is not a return point.
\item \textbf{Controllable exhaustion at failure.} A safe non-goal state with no uncontrollable has a cursor over all controllable candidates. Expansion consumes it to the next nonempty bucket or exhaustion. If a nonwinner has candidates remaining, it is in $\mathsf{Retry}$. Each retry consumes further candidates. The failure return is reached only after both queues and the propagation worklist are empty. Hence every such state in $V\setminus W$ has exhausted its cursor. Removing an entry when its state becomes winning cannot remove a losing candidate from this argument.
\item \textbf{Monotone winning evidence.} A promoted state either has all uncontrollable outcomes at smaller retained ranks, or has a complete selected controllable bucket at smaller ranks and no uncontrollable. All those buckets are exact; growing $\mathcal D$ adds other controllable choices rather than new outcomes to an already stored choice. The policy disables those additional controllables. Existing rank-decreasing witness edges consequently remain a valid certificate. Goals are recognized using modeled uncontrollable enabledness before terminalization, so an unqueried uncontrollable cannot invalidate rank zero.
\item \textbf{Open-state coverage.} Each newly discovered safe non-goal is inserted once in $\mathsf{Open}$. It is removed only for first expansion, when it enters $X$. Thus $V\cap\Safe\setminus(X\cup\Goal)\subseteq Open$. Unsafe states and goals are discovered terminals and are handled separately.
\end{enumerate}
The initial root insertion establishes these properties. Expansion, successor discovery, candidate consumption, and promotion preserve them as shown above; no other operation changes their relevant sets.

